\documentclass{article}

\usepackage{graphicx} % Required for inserting images
\usepackage{xcolor}
\usepackage{amssymb}
\usepackage{amsmath} % for \DeclareMathOperator
\usepackage{amsthm}
\usepackage{float}
\usepackage{booktabs}
\usepackage{cite}

\floatstyle{ruled}
\newfloat{algorithm}{tbp}{loa}
\floatname{algorithm}{Algorithm}

\newtheorem{theorem}{Theorem}
\newtheorem{lemma}{Lemma}
\newtheorem{definition}{Definition}
\newtheorem{corollary}{Corollary}
\newtheorem{remark}{Remark}

\DeclareMathOperator{\lcm}{lcm}
%\DeclareMathOperator{\max}{max}

\newcommand{\ceil}[1]{\left\lceil #1 \right\rceil}
\newcommand{\Comment}[1]{\hfill\textnormal{\textit{// #1}}}


\title{LoRa AOI}
%\author{nirajcse08 }
\date{ }
\newcommand{\nkt}[1]{\textcolor{blue}{#1}}

\newcommand{\tbd}{{\color{blue}TBD}}

%% Notations

%index
\newcommand{\taskInd}{\ensuremath{i}} %index for tasks
\newcommand{\taskIndTwo}{\ensuremath{k}} %another index for tasks
\newcommand{\jobInd}{\ensuremath{x}} %index for jobs
\newcommand{\procInd}{\ensuremath{j}} %index for processors
\newcommand{\freqInd}{\ensuremath{q}} %freq index
\newcommand{\tempInd}{\ensuremath{z}} %anywhere

%input
\newcommand{\task}{\ensuremath{\tau}}   %task
\newcommand{\period}{\ensuremath{p}}   %period of task/deadline of job
\newcommand{\execTime}{\ensuremath{e}}   %execution time of task
\newcommand{\proc}{\ensuremath{P}}   %processor
\newcommand{\freq}{\ensuremath{f}}   %frequency
\newcommand{\execTimeAttack}{\ensuremath{e}^{attack}}   %execution time of task
\newcommand{\util}{\ensuremath{u}}   % utility

%intermediate
\newcommand{\prior}{\ensuremath{\rho}}   %priority of a task
\newcommand{\finTime}{\ensuremath{\delta}}   %letter for finishing time 
\newcommand{\hyper}{\ensuremath{h}}   %letter for denoting heaviness of a job
\newcommand{\fmax}{\ensuremath{\freq}^{max}}   %max frequency
\newcommand{\energy}{\ensuremath{E}} % letter for energy consumption

%set
\newcommand{\taskSet}{\ensuremath{\mathcal{T}}} %task set
\newcommand{\critTaskSet}{\ensuremath{\mathcal{T}^{crit}}} % safety critical task set
\newcommand{\nonCritTaskSet}{\ensuremath{\mathcal{T}^{ncrit}}} %non-safety critical task set
\newcommand{\procSet}{\ensuremath{\mathcal{P}}} %proc set
\newcommand{\freqSet}{\ensuremath{\mathcal{F}}} %freq set
\newcommand{\hpTasks}{\ensuremath{\mathcal{H}}}   %set of hp jobs of Ji
\newcommand{\lpTasks}{\ensuremath{\mathcal{L}}}   %letter for set of lp jobs of Ji
\newcommand{\selectedJobs}{\ensuremath{\mathcal{S}}} % selected set of jobs
\newcommand{\jobSet}{\ensuremath{\mathcal{J}}} % jobs assigned to a processor
\newcommand{\arrTimeSet}{\ensuremath{\mathcal{A}}} % all arrivals of jobs
\newcommand{\deadlineSet}{\ensuremath{\mathcal{D}}} % deadline for all jobs
\newcommand{\jobPairSet}{\ensuremath{\mathcal{B}}} % pair of jobs
\newcommand{\victimSet}{\ensuremath{\mathcal{V}}} % set of victim tasks
\newcommand{\nonVictimSet}{\ensuremath{\mathcal{NV}}} % non victim set
\newcommand{\safetyCrit}{\ensuremath{\mathcal{SC}}} % set of safety crit tasks
\newcommand{\nonSafetyCrit}{\ensuremath{\mathcal{NSC}}} % non safety crit tasks


%count
\newcommand{\noTasks}{\ensuremath{N^{tsk}}} %number of tasks
\newcommand{\noCritTasks}{\ensuremath{N^{ctsk}}} %number of safety crit tasks
\newcommand{\noNonCritTasks}{\ensuremath{N^{nctsk}}} %number of non-sfety critical tasks
\newcommand{\noJobs}{\ensuremath{N^{job}}} %number of jobs
\newcommand{\noProc}{\ensuremath{N^{proc}}} %number of processors
\newcommand{\noFreq}{\ensuremath{N^{freq}}} %number of freq

%formulation
\newcommand{\bvar}{\ensuremath{X}} %decision variable 
\newcommand{\obj}{\ensuremath{\Psi}} % objective 

\newcommand{\curSoln}{\ensuremath{S}^{cur}} % current solution
\newcommand{\nextSoln}{\ensuremath{S}^{next}} % next solution
\newcommand{\bestSoln}{\ensuremath{S}^{best}} % best solution
\newcommand{\tempSoln}{\ensuremath{S}^{temp}} % temp solution
\newcommand{\selectedSoln}{\ensuremath{S}^{sel}} % temp solution

\newcommand{\curSolnCost}{\ensuremath{C}^{cur}} % current solution cost
\newcommand{\nextSolnCost}{\ensuremath{C}^{next}} % next solution cost
\newcommand{\bestSolnCost}{\ensuremath{C}^{best}} % best solution cost
\newcommand{\tempSolnCost}{\ensuremath{C}^{temp}} % temp solution cost
\newcommand{\selectedSolnCost}{\ensuremath{C}^{sel}} % temp solution cost

\newcommand{\probBadSoln}{\ensuremath{\beta}} % prob of accepting a bad solution
\newcommand{\probFreqDec}{\ensuremath{\gamma}} % prob of frequency decreaase
\newcommand{\probScalingFactor}{\ensuremath{q}} % prob of frequency decreaase

\newcommand{\dbf}{\ensuremath{DBF}}

%BV
\newcommand{\dvar}{\ensuremath{V}} %letter denoting binary variable for task
\newcommand{\dvarAux}{\ensuremath{Y}} %letter denoting aux binary variable

\newcommand{\releaseTime}{\ensuremath{r}} %release time of a job = (j-1)period
\newcommand{\deadline}{\ensuremath{d}} %deadline of a job = (j)period
\newcommand{\effExecTime}{\ensuremath{e^{eff}}} %letter denoting effective exec time of a job with specific freq and optional segment
\newcommand{\exec}{\ensuremath{e}} %letter denoting exec time
%\newcommand{\execOptional}{\ensuremath{e^{O}}} %letter denoting optional execution time
\newcommand{\optionalSeg}{\ensuremath{G}} %letter denoting optional segment
\newcommand{\totalEnergy}{\ensuremath{E^{total}}} %total denoting energy consumption
\newcommand{\utilTask}{\ensuremath{u}} %letter denoting utility of a task, utility gained per unit execution of optional segment
\newcommand{\execStartTime}{\ensuremath{t^{beg}}} %letter denoting execution start time
\newcommand{\execFinishTime}{\ensuremath{t^{fin}}} %letter denoting execution finish time
\newcommand{\energyBudget}{\ensuremath{\mathcal{B}}} %letter denoting energy budget
\newcommand{\numPostions}{\ensuremath{L}} %letter denoting positions
\newcommand{\numAlloc}{\ensuremath{A}} %letter denoting allocations on a processor
\newcommand{\timeSlack}{\ensuremath{\Delta^t}} % time slack
\newcommand{\energySlack}{\ensuremath{\Delta^e}} % energy slack
\newcommand{\selectedJob}{\ensuremath{J^*}} % a job 
\newcommand{\utilGainedJob}{\ensuremath{U^{job}}} % utility gained by a job
\newcommand{\utilGainedTask}{\ensuremath{U^{tsk}}} % utility gained by a task
\newcommand{\utilGainedTotal}{\ensuremath{U^{total}}} % total utility gained
\newcommand{\jobPair}{\ensuremath{p}} % a job pair


\title{\textbf{A Combinatorial Approximation Framework for Real-Time Packet Scheduling in LoRaWAN Networks}}
\author{BTP Project Report}
\date{\today}

\begin{document}
\maketitle

\section{System Model and Problem Statement}
\label{sec:model}

\subsection{Network and Hardware Model}
Consider a star-of-stars LoRaWAN deployment comprising a set of distributed end-devices communicating with a central multi-channel gateway. The network operates under a cyclic hypercycle structure of duration $H$. 

The physical transmission medium is characterized by:
\begin{itemize}
    \item A set of $N_{\text{ch}}$ orthogonal carrier frequencies (channels) $\mathcal{C} = \{\text{CH}_0, \text{CH}_1, \dots, \text{CH}_{N_{\text{ch}}-1}\}$.
    \item A set of $N_{\text{sf}}$ quasi-orthogonal spreading factors $\mathcal{S} = \{\text{SF}_7, \text{SF}_8, \dots, \text{SF}_{12}\}$.
\end{itemize}

A physical communication resource is uniquely defined by a tuple $\langle \text{SF}, \text{CH} \rangle$. We model the gateway receiver capabilities as a set of $m = N_{\text{ch}} \times N_{\text{sf}}$ discrete, unrelated parallel machines:
\begin{equation}
\mathcal{M} = \{M_0, M_1, \dots, M_{m-1}\}
\end{equation}
For an 8-channel gateway operating across spreading factors SF7--SF12, the system comprises $m = 8 \times 6 = 48$ logical machines.

Hardware constraints at standard gateway digital baseband processors (e.g., Semtech SX1301/SX1302/SX1303) restrict the maximum number of concurrent demodulation paths across all channels and spreading factors to at most \textbf{8 simultaneous packets}:
\begin{equation}
N_{\text{demod}}^{\max} = 8
\end{equation}
Overlapping transmissions exceeding this threshold induce baseband receiver channel exhaustion, triggering buffer overflow, DSP pipeline saturation, and unrecoverable packet drops.

\subsection{Traffic and Packet Model}
Let $\mathcal{J} = \{J_0, J_1, \dots, J_{n-1}\}$ denote the set of $n$ real-time packets to be scheduled within the hypercycle $[0, H]$. The traffic is periodic with harmonic periods: each packet stream $k$ possesses a period $T_k$ such that $T_k \mid H$, generating periodic instances with release times $r_i \in [0, H)$.

Each packet $J_i \in \mathcal{J}$ is characterized by:
\begin{enumerate}
    \item A release time $r_i \ge 0$, representing the earliest instant the end-device can initiate transmission.
    \item An absolute deadline $D_i \in (r_i, H]$, clamped to the hypercycle boundary:
    \begin{equation}
    D_i \leftarrow \min(D_i, H)
    \end{equation}
    \item A criticality or priority level $\Pi_i \in \mathbb{R}^+$, where higher values designate packets of greater operational importance.
    \item A payload size $L_i$ in bytes.
\end{enumerate}

\subsection{Time-on-Air (ToA) Matrix}
The transmission duration $p_{ij}$ of packet $J_i$ on resource $M_j = \langle \text{SF}_j, \text{CH}_j \rangle$ is governed by the Semtech SX1276 physical layer specifications:
\begin{equation}
T_{\text{sym}} = \frac{2^{\text{SF}_j}}{\text{BW}} \quad (\text{seconds})
\end{equation}
The preamble duration is $T_{\text{preamble}} = (N_{\text{preamble}} + 4.25) \cdot T_{\text{sym}}$. The payload symbol count is given by:
\begin{equation}
N_{\text{payload}} = 8 + \max\left( \left\lceil \frac{8 L_i - 4\text{SF}_j + 28 + \text{CRC} - 20\text{IH}}{4(\text{SF}_j - 2\text{DE})} \right\rceil (\text{CR} + 4), 0 \right)
\end{equation}
where $\text{CRC} = 16$ bits, $\text{IH} = 0$ (explicit header enabled), $\text{CR} = 1$ (coding rate 4/5), and $\text{DE} = 1$ if $T_{\text{sym}} > 16\text{ ms}$ (Low Data Rate Optimization), and 0 otherwise.

The entries of the $n \times m$ Time-on-Air (ToA) matrix $P = [p_{ij}]$ are discretized into integer milliseconds:
\begin{equation}
p_{ij} = \begin{cases}
\lceil (T_{\text{preamble}} + N_{\text{payload}} \cdot T_{\text{sym}}) \times 1000 \rceil, & \text{if packet } J_i \text{ is eligible on } M_j \\
\infty, & \text{otherwise}
\end{cases}
\end{equation}
We define the maximum processing time and the number of admissible edges as:
\begin{equation}
U = \max_{i \in \mathcal{J}, j \in \mathcal{M}} \{ p_{ij} \mid p_{ij} < \infty \}, \quad A = |\{(i, j) \mid p_{ij} < \infty\}| \le n \cdot m
\end{equation}

\subsection{LoRaWAN Class A Energy Consumption Model}
\label{subsec:energy_model}
Energy efficiency is paramount for battery-powered LoRaWAN IoT end-devices. Following the established analytical frameworks of Casals et al.~\cite{casals2017modeling} and Bouguera et al.~\cite{bouguera2018energy}, the energy consumption per transmission cycle $E_{\text{cycle}}$ encompasses all operational radio and microcontroller states:
\begin{equation}
E_{\text{cycle}} = V_{\text{dd}} \cdot \sum_{s \in \mathcal{S}_{\text{states}}} \left( I_s \cdot t_s \right)
\end{equation}
where $\mathcal{S}_{\text{states}} = \{\text{MCU}, \text{Sensor}, \text{TX}, \text{Wait1}, \text{RX1}, \text{Wait2}, \text{RX2}, \text{Sleep}\}$. Specifically:
\begin{equation}
E_{\text{tx}}(i, j) = V_{\text{dd}} \cdot I_{\text{tx}} \cdot p_{ij}
\end{equation}
Because the symbol duration $T_{\text{sym}} = 2^{\text{SF}} / \text{BW}$ grows exponentially with Spreading Factor, Time-on-Air and transmit energy scale sharply. For instance, for a $10$-byte application payload ($23$ bytes total PHY payload with LoRaWAN MAC overhead), transmitting at $\text{SF}_{12}$ consumes $233.7\text{ mJ}$ compared to only $27.4\text{ mJ}$ at $\text{SF}_7$ ($\approx 8.5\times$ energy inflation).

Consequently, our scheduling objective is two-fold:
\begin{enumerate}
    \item \textbf{Makespan Guarantee:} Bound channel congestion such that $C_{\max}^{\text{mach}}(\alpha) \le 2 \cdot \text{OPT}(P)$.
    \item \textbf{Energy Minimization:} Minimize total fleet radio energy $\sum_{i} E_{\text{cycle}}(i, \alpha(i))$ by aggressively downshifting packets to lower Spreading Factors without violating the makespan upper bound.
\end{enumerate}

\subsection{Problem Formulation}
An assignment of packets to resources is defined by a mapping function:
\begin{equation}
\alpha: \mathcal{J} \to \mathcal{M}
\end{equation}
The total workload on resource $M_j$ under assignment $\alpha$ is the cumulative transmission duration of all packets assigned to $j$:
\begin{equation}
\delta_j(P, \alpha) = \sum_{i \in \mathcal{J}, \alpha(i) = j} p_{ij}
\end{equation}

To avoid ambiguity between scheduling metrics, we explicitly distinguish:
\begin{definition}[Machine Load Makespan]
The maximum channel workload across all parallel resources:
\begin{equation}
C_{\max}^{\text{mach}}(\alpha) = \max_{j \in \mathcal{M}} \delta_j(P, \alpha), \quad \text{OPT}(P) = \min_{\alpha} C_{\max}^{\text{mach}}(\alpha)
\end{equation}
\end{definition}

\begin{definition}[Timeline Calendar Span]
The latest chronological packet completion time along the hypercycle timeline:
\begin{equation}
C_{\max}^{\text{time}} = \max_{i \in \mathcal{J}_{\text{sched}}} f_i \le H
\end{equation}
\end{definition}

The combinatorial 2-approximation theorem applies specifically to minimizing the machine load makespan $C_{\max}^{\text{mach}}(\alpha)$. In addition to load balancing, each resource $M_j$ represents a single non-preemptive communication channel. To satisfy real-time constraints, the set of packets $\mathcal{J}_j = \{J_i \in \mathcal{J} \mid \alpha(i) = j\}$ assigned to resource $M_j$ must be chronologically schedulable without deadline violations and without violating the gateway's concurrent demodulation threshold ($N_{\text{demod}}^{\max} = 8$).

\section{Combinatorial 2-Approximation on Unrelated Machines}
\label{sec:approx_theory}

We apply the purely combinatorial polynomial-time 2-approximation algorithm of Gairing, Monien, and Woclaw \cite{gairing2007faster} for scheduling on unrelated parallel machines ($R || C_{\max}$).

\subsection{Generalized Unsplittable Flow Formulation}
Given a candidate makespan $w \in \mathbb{Z}^+$, we construct a bipartite directed network with node set $V = \mathcal{J} \cup \mathcal{M}$, where job nodes $\mathcal{J}$ represent packets and machine nodes $\mathcal{M}$ represent resources.
Each job node $i \in \mathcal{J}$ possesses a supply of 1 unit. An edge from job $i$ to machine $j$ exists if and only if $p_{ij} \le w$, with unit capacity and multiplicative gain $\mu_{ij} = p_{ij}$. 

A generalized unsplittable flow routes the entire 1 unit of supply from each job node $i$ to exactly one machine $j$, creating an accumulated load $\delta_j(P, \alpha)$ on machine $j$. An assignment $\alpha$ is termed \emph{$w$-feasible} if $p_{i, \alpha(i)} \le w$ for all $i \in \mathcal{J}$.

\subsection{Residual Network and Machine Partitioning}
For a current $w$-feasible assignment $\alpha$, we define the directed residual graph $G_\alpha(w) = (V, E_\alpha(w))$ with edge set $E_\alpha(w) = E_\alpha^1 \cup E_\alpha^2$:
\begin{itemize}
    \item \textbf{Backward Edges ($E_\alpha^1$):} For each assigned job $\alpha(i) = j$, add directed edge $(j, i)$ of capacity 1 and gain $\mu_{ji} = 1/p_{ij}$. This represents removing job $i$ from machine $j$.
    \item \textbf{Forward Edges ($E_\alpha^2$):} For each unassigned pair where $\alpha(i) \neq j$ and $p_{ij} \le w$, add directed edge $(i, j)$ of capacity 1 and gain $\mu_{ij} = p_{ij}$. This represents reassigning job $i$ to machine $j$.
\end{itemize}

Given the current loads $\delta_j(P, \alpha)$, the machines are partitioned into three mutually disjoint sets:
\begin{definition}[Machine Partitioning]
\label{def:partitions}
\begin{align}
\mathcal{M}^-(\alpha) &= \{ j \in \mathcal{M} \mid \delta_j(P, \alpha) \le w \} \quad && \text{(Underloaded / Sinks)} \\
\mathcal{M}^0(\alpha) &= \{ j \in \mathcal{M} \mid w + 1 \le \delta_j(P, \alpha) \le 2w \} \quad && \text{(Medium / Transit)} \\
\mathcal{M}^+(\alpha) &= \{ j \in \mathcal{M} \mid \delta_j(P, \alpha) \ge 2w + 1 \} \quad && \text{(Overloaded / Sources)}
\end{align}
\end{definition}

If $\mathcal{M}^+(\alpha) = \emptyset$, every machine satisfies $\delta_j(P, \alpha) \le 2w$, which yields a valid 2-approximation provided $\text{OPT}(P) \ge w$.

\subsection{Gain Rounding and Truemper's Logarithmic Transformation}
To apply minimum-cost flow techniques without introducing numerical instability, edge gains are geometrically rounded down to integer powers of a base $b > 1$:
\begin{equation}
b = 1 + \frac{1}{m}
\end{equation}
For each edge $(i, j)$, the cost is defined by the discrete logarithm:
\begin{equation}
c_{ij} = \lfloor \log_b(p_{ij}) \rfloor \in \mathbb{Z}
\end{equation}
The rounded gain is $\gamma_{ij} = b^{c_{ij}}$. For reverse edges $(j, i)$, antisymmetry is enforced: $c_{ji} = -c_{ij}$ and $\gamma_{ji} = 1/\gamma_{ij} = b^{-c_{ij}}$.

Crucially, the rounding ensures:
\begin{equation}
\label{eq:rounding_bound}
\gamma_{ij} \le p_{ij} \le b \cdot \gamma_{ij}
\end{equation}
The relative error introduced per edge is strictly bounded by $b - 1 = 1/m$.

\subsection{Node Potentials and Reduced Cost Optimality}
A node potential function $\pi: V \to \mathbb{R}$ is maintained across all nodes. For any directed edge $(u, v)$, the reduced cost is defined as:
\begin{equation}
c_{uv}^\pi = c_{uv} - \pi(u) + \pi(v)
\end{equation}
The Primal-Dual algorithm preserves the reduced cost optimality condition:
\begin{equation}
c_{uv}^\pi \ge 0 \quad \forall (u, v) \in E_\alpha(w)
\end{equation}
This condition guarantees that no negative-cost cycles exist in the residual network with respect to edge costs $c$.

\section{Algorithmic Mechanics and Theoretical Proofs}
\label{sec:algorithms_proofs}

The approximation framework operates via two nested algorithmic procedures: the inner loop, \emph{Unsplittable-Blocking-Flow}, and the outer loop, \emph{Unsplittable-Truemper}.

\subsection{Inner Loop: Unsplittable-Blocking-Flow}
The inner procedure pushes discrete jobs along alternating paths from overloaded machines $\mathcal{M}^+$ toward underloaded machines $\mathcal{M}^-$ on a shortest-path level graph.

Let $G'_\alpha(w)$ be an admissible subgraph consisting of edges with zero reduced cost ($c_{uv}^\pi = 0$).
A height function $h: V \to \mathbb{N} \cup \{\infty\}$ is computed, where $h(v)$ represents the shortest hop distance from node $v$ to any sink in $\mathcal{M}^-$ in $G'_\alpha(w)$.
An edge $(u, v)$ is \emph{admissible} if:
\begin{equation}
h(u) = h(v) + 1
\end{equation}

\begin{definition}[Helpful Machine and Helpful Path]
A machine $j \in \mathcal{M}^+$ is a \textbf{helpful machine} if it can reassign its first admissible job $s(j)$ while keeping its remaining load at or above $w+1$:
\begin{equation}
\delta_j(P, \alpha) \ge w + 1 + p_{s(j), j}
\end{equation}
A \textbf{helpful path} is an alternating sequence of nodes $v_0, v_1, \dots, v_r$ where even indices $v_{2i} \in \mathcal{M}$ are machines and odd indices $v_{2i+1} \in \mathcal{J}$ are jobs, such that:
\begin{enumerate}
    \item $v_0 \in \mathcal{M}^+$ is a helpful machine of minimum height.
    \item Every edge $(v_i, v_{i+1})$ is admissible: $h(v_i) = h(v_{i+1}) + 1$.
    \item For all intermediate machines ($0 < i < r/2$):
    \begin{equation}
    w + 1 \le \delta_{v_{2i}} + p_{s(v_{2i-2}), v_{2i}} - p_{s(v_{2i}), v_{2i}} \le 2w
    \end{equation}
    \item For the terminal machine $v_r \in \mathcal{M}^-$:
    \begin{equation}
    \delta_{v_r} + p_{s(v_{r-2}), v_r} \le 2w
    \end{equation}
\end{enumerate}
\end{definition}

Reassigning jobs along a helpful path shifts one job from each machine to the next. The load of $v_0$ decreases but remains $\ge w+1$; intermediate machines remain strictly within $\mathcal{M}^0$; and the load of $v_r$ increases but never exceeds $2w$.

\begin{algorithm}[t]
\caption{Unsplittable-Blocking-Flow$(\alpha, G'_\alpha(w), P, w)$}
\label{alg:ubf}
\small
\textbf{Input:} Assignment $\alpha$, admissible graph $G'_\alpha(w)$, processing matrix $P$, parameter $w$\\
\textbf{Output:} Updated assignment $\beta$
\begin{enumerate}\setlength{\itemsep}{1pt}\setlength{\parskip}{0pt}
    \item Compute height labels $h(v)$ for all $v \in V$ as shortest hop distance to $\mathcal{M}^-$ in $G'_\alpha(w)$ via backward BFS.
    \item \textbf{while} $\mathcal{M}^- \neq \emptyset$ \textbf{and} $\exists j \in \mathcal{M}^+$ with $h(j) < \infty$ \textbf{do}
    \item \quad Identify helpful machines: $\mathcal{H} \leftarrow \{j \in \mathcal{M}^+ \mid h(j) < \infty, \, \delta_j \ge w + 1 + p_{s(j), j}\}$.
    \item \quad \textbf{if} $\mathcal{H} = \emptyset$ \textbf{then} \textbf{break}
    \item \quad Select a helpful machine $v_0 \in \mathcal{H}$ of minimum height.
    \item \quad Construct helpful path induced by $v_0$.
    \item \quad Push jobs along the helpful path: update $\alpha$, machine loads $\delta$, and admissible graph edges.
    \item \quad Recompute height labels $h$ via BFS.
    \item \textbf{end while}
    \item \textbf{return} $\alpha$
\end{enumerate}
\end{algorithm}

\begin{lemma}[Load Invariance \cite{gairing2007faster}]
\label{lem:load_invariance}
Let $\beta$ be the assignment returned by Algorithm~\ref{alg:ubf} ($\text{Unsplittable-Blocking-Flow}$). Then:
\begin{enumerate}
    \item If $j \in \mathcal{M}^-(\alpha)$, then $\delta_j(P, \\beta) \ge \delta_j(P, \alpha)$ (sink loads never decrease).
    \item If $j \in \mathcal{M}^0(\\alpha)$, then $w + 1 \le \delta_j(P, \beta) \le 2w$ (medium machines stay medium).
    \item If $j \in \mathcal{M}^+(\\alpha)$, then $\delta_j(P, \beta) \le \delta_j(P, \alpha)$ (source loads never increase).
\end{enumerate}
\end{lemma}

\begin{theorem}[Complexity of Inner Loop \cite{gairing2007faster}]
\label{thm:ubf_complexity}
The inner loop terminates in at most $m$ phases. Since each phase requires a Breadth-First Search (BFS) of complexity $O(A)$, the total running time of $\text{Unsplittable-Blocking-Flow}$ is $O(mA)$. Upon termination, no path exists from $\mathcal{M}^+(\beta)$ to $\mathcal{M}^-(\beta)$ in $G'_\beta(w)$.
\end{theorem}

\subsection{Outer Loop: Unsplittable-Truemper}
The outer loop maintains node potentials $\pi(v)$ and iteratively constructs the admissible shortest-path graph for the inner loop using reduced-cost Dijkstra on the reversed residual graph.

\begin{algorithm}[t]
\caption{Unsplittable-Truemper$(\alpha, P, C, w)$}
\label{alg:truemper}
\small
\textbf{Input:} Initial $w$-feasible assignment $\alpha$, matrices $P$ and $C$, parameter $w$\\
\textbf{Output:} Final assignment $\beta$, or $\emptyset$ if $\text{OPT}(P) > w$
\begin{enumerate}\setlength{\itemsep}{1pt}\setlength{\parskip}{0pt}
    \item Initialize potentials: $\pi(v) \leftarrow 0$ for all $v \in \mathcal{M}$; $\pi(i) \leftarrow C_{i, \alpha(i)}$ for all $i \in \mathcal{J}$.
    \item \textbf{while} $\mathcal{M}^+ \neq \emptyset$ \textbf{do}
    \item \quad \textbf{if} $\max_{j \in \mathcal{M}^+} \pi(j) \ge \log_b(m)$ \textbf{then} \textbf{return} $\emptyset$ \Comment{Potential ceiling reached}
    \item \quad Compute shortest distances $d(v)$ from all nodes to sinks $\mathcal{M}^-$ in reversed residual graph via Dijkstra using reduced costs $c_{uv}^\pi$.
    \item \quad \textbf{if} $\forall j \in \mathcal{M}^+: d(j) = \infty$ \textbf{then} \textbf{return} $\emptyset$ \Comment{No path from $\mathcal{M}^+$ to $\mathcal{M}^-$}
    \item \quad Update potentials: $\pi(v) \leftarrow \pi(v) + d(v)$ for all $v \in V$ with $d(v) < \infty$.
    \item \quad Construct admissible graph $G'_\alpha(w)$ containing edges with $c_{uv}^\pi = 0$.
    \item \quad $\alpha_{\text{prev}} \leftarrow \alpha$.
    \item \quad $\alpha \leftarrow \text{Unsplittable-Blocking-Flow}(\alpha, G'_\alpha(w), P, w)$.
    \item \quad \textbf{if} $\alpha = \alpha_{\text{prev}}$ and all $d(j) = 0$ on reachable sources \textbf{then} \textbf{return} $\emptyset$.
    \item \textbf{end while}
    \item \textbf{return} $\alpha$
\end{enumerate}
\end{algorithm}

\begin{theorem}[Lower Bound Certificate \cite{gairing2007faster}]
\label{thm:lower_bound}
If Algorithm~\ref{alg:truemper} terminates with $\mathcal{M}^+ \neq \emptyset$ due to the potential ceiling $\max_{j \in \mathcal{M}^+} \pi(j) \ge \log_b(m)$ or because no path connects $\mathcal{M}^+$ to $\mathcal{M}^-$, then no fractional or unsplittable assignment of makespan $\le w$ exists; hence $\text{OPT}(P) \ge w + 1$.
\end{theorem}

\subsection{Binary Search and 2-Approximation Theorem}
\begin{theorem}
\label{thm:overall_2approx}
Algorithm~\ref{alg:truemper} combined with binary search computes an assignment $\alpha$ satisfying:
\begin{equation}
C_{\max}^{\text{mach}}(\alpha) \le 2 \cdot \text{OPT}(P)
\end{equation}
in total time $O(m^2 A \log m \log(w_{\max} - w_{\min}))$.
\end{theorem}
\begin{proof}
We perform binary search over $w \in [w_{\min}, w_{\max}]$, where $w_{\min} = \max_{i} \min_j p_{ij}$ and $w_{\max} = 2 \lceil \sum_i \min_j p_{ij} / m \rceil + 2 \max_{i,j} p_{ij}$. We locate the critical threshold $w^*$ such that Algorithm~\ref{alg:truemper} with parameter $w^* - 1$ terminates with $\mathcal{M}^+ \neq \emptyset$ (proving $\text{OPT}(P) \ge w^*$ by Theorem~\ref{thm:lower_bound}), while with parameter $w^*$ it terminates with $\mathcal{M}^+ = \emptyset$.

When $\mathcal{M}^+ = \emptyset$, every machine satisfies $\delta_j(P, \alpha) \le 2w^*$.
Thus:
\begin{equation}
C_{\max}^{\text{mach}}(\alpha) \le 2w^* \le 2 \cdot \text{OPT}(P)
\end{equation}
The binary search requires at most $O(\log(w_{\max} - w_{\min}))$ calls. With each call costing $O(m^2 A \log m)$, the overall running time is $O(m^2 A \log m \log(w_{\max} - w_{\min}))$.
\end{proof}

\section{LoRaWAN Hypercycle Scheduling with Hardware Demodulation Constraints}
\label{sec:lorawan_scheduler}

\subsection{Gateway Hardware Demodulation Constraint}
\label{subsec:hardware_limit}
In standard LoRaWAN deployments, the central gateway is equipped with multi-channel digital baseband processors (such as Semtech SX1301/SX1302/SX1303). Although the gateway listens across all spreading factors $\mathcal{S} = \{\text{SF}_7, \dots, \text{SF}_{12}\}$ defining $m = N_{\text{ch}} \times N_{\text{sf}} = 48$ logical communication pipes $\mathcal{M}$, the physical modem architecture possesses a finite pool of demodulation paths, strictly capping the maximum number of simultaneous packet receptions across all channels to at most \textbf{8 concurrent demodulations}:
\begin{equation}
N_{\text{demod}}^{\max} = 8
\end{equation}

When more than 8 packets overlap in time across all frequency channels and spreading factors, the baseband modem experiences receiver channel exhaustion, causing packet drops due to hardware buffer overflow and DSP pipeline saturation. Therefore, any physically realizable schedule must satisfy two complementary conditions:
\begin{enumerate}
    \item \textbf{Single-Channel Non-Preemption:} On any individual resource $M_j = \langle \text{SF}_j, \text{CH}_j \rangle$, transmissions cannot overlap in time.
    \item \textbf{Global Modem Demodulation Limit:} Across all $m$ resources, the total number of simultaneous active transmissions at any continuous instant $t \in [0, H]$ must not exceed 8.
\end{enumerate}

Let $[s_i, f_i)$ represent the scheduled transmission interval of packet $J_i$ on its assigned resource $M_j = \alpha(i)$, where $f_i = s_i + p_{ij}$. The global gateway demodulation occupancy at instant $t$ is:
\begin{equation}
\Omega(t) = \\sum_{i \in \mathcal{J}_{\text{sched}}} \mathbb{I}(s_i \le t < f_i)
\end{equation}
The global hardware constraint requires:
\begin{equation}
\label{eq:demod_constraint}
\Omega(t) \le 8 \quad \forall t \in [0, H]
\end{equation}

\subsection{Per-Channel Earliest Deadline First (EDF) Timeline}
\label{subsec:edf_timeline}
For each individual machine $M_j \in \mathcal{M}$, the assigned packets $\mathcal{J}_j = \{J_i \in \mathcal{J}_{\text{cand}} \mid \alpha(i) = j\}$ are sorted in ascending order of deadlines (Earliest Deadline First):
\begin{equation}
D_{(1)} \le D_{(2)} \le \dots \le D_{(k)}
\end{equation}

To maintain $\Omega(t) \le 8$ without expensive continuous sweeps, we implement a Segment Tree with lazy propagation over $[0, H]$. 
A min-heap priority queue coordinates multi-channel progression:
\begin{align}
s_i &= \max\left(\text{cursor}_j, r_i\right) \\
f_i &= s_i + p_{ij}
\end{align}
The demodulation peak is checked in $O(\log H)$ time via $\text{range\_max}(s_i, f_i - 1)$. If $\text{peak} \ge 8$ or $f_i > D_i$, the packet is flagged as an offender. Otherwise, the occupancy is updated via $\text{range\_add}(s_i, f_i - 1, +1)$ in $O(\log H)$ time, and $\text{cursor}_j \leftarrow f_i$.

\subsection{Phase 1.5: Energy-Optimal Capacity-Constrained Downshifting}
\label{subsec:phase1_5_energy}
While Phase 1 produces a balanced assignment $\alpha$ satisfying $C_{\max}^{\text{mach}}(\alpha) \le 2 \cdot \text{OPT}(P)$, it optimizes makespan without explicit regard to radio energy. Because energy is heavily penalized on higher Spreading Factors ($E_{i, \text{SF12}} \gg E_{i, \text{SF7}}$), many packets can be migrated to lower Spreading Factors without exceeding the machine capacity budget.

We formulate the post-Phase 1 energy optimization problem as follows:
\begin{align}
\min_{\alpha^*} \quad & \sum_{i \in \mathcal{J}_{\text{cand}}} E_{\text{cycle}}(i, \alpha^*(i)) \label{opt:energy_obj} \\
\text{subject to} \quad & \sum_{i: \alpha^*(i) = j} p_{ij} \le 2 w^* \quad \forall j \in \mathcal{M} \label{opt:energy_cap} \\
& p_{i, \alpha^*(i)} \le w^* \quad \forall i \in \mathcal{J}_{\text{cand}} \label{opt:energy_feas}
\end{align}

\begin{theorem}[Preservation of the 2-Approximation Upper Bound under Phase 1.5]
\label{thm:energy_preserves_bound}
Let $\alpha^*$ be any assignment produced by Phase~1.5 satisfying constraint~\eqref{opt:energy_cap}. Then:
\begin{equation}
C_{\max}^{\text{mach}}(\alpha^*) \le 2 \cdot \text{OPT}(P)
\end{equation}
Furthermore, the total radio energy is monotonically non-increasing:
\begin{equation}
\sum_{i} E_{\text{cycle}}(i, \alpha^*(i)) \le \sum_{i} E_{\text{cycle}}(i, \alpha(i))
\end{equation}
\end{theorem}
\begin{proof}
By Theorem~\ref{thm:lower_bound}, the critical threshold found by binary search provides a lower bound on the optimum: $\text{OPT}(P) \ge w^*$. Constraint~\eqref{opt:energy_cap} enforces that no machine $j \in \mathcal{M}$ incurs total processing time exceeding $2 w^*$. Therefore:
\begin{equation}
C_{\max}^{\text{mach}}(\alpha^*) = \max_{j \in \mathcal{M}} \delta_j(P, \alpha^*) \le 2 w^* \le 2 \cdot \text{OPT}(P)
\end{equation}
Because the initial assignment $\alpha$ from Phase 1 already satisfies $\delta_j(P, \alpha) \le 2 w^*$, $\alpha$ is a feasible solution to the constrained optimization problem. Since Phase 1.5 accepts only transitions that yield non-negative energy reduction ($\Delta E \ge 0$), the total energy cannot increase.
\end{proof}

\subsection{Complete Main Algorithm Pseudocode}
\label{subsec:main_algorithm}
The complete end-to-end scheduling framework is formalized in Algorithm~\ref{alg:lorawan_scheduler}. It integrates makespan minimization via $\text{Unsplittable-Truemper}$, energy-optimal capacity-constrained downshifting (Phase 1.5), per-channel EDF timeline construction, Segment-Tree-driven 8-demodulation verification, and priority-conscious packet dropping.

\begin{algorithm}[t!]
\caption{LoRaWAN-Hypercycle-Scheduler$(\mathcal{J}, \mathcal{M}, P, E, H, N_{\text{demod}}^{\max})$}
\label{alg:lorawan_scheduler}
\small
\textbf{Input:} Packet set $\mathcal{J}$ with release times $r_i$, deadlines $D_i$, priorities $\Pi_i$; Resource set $\mathcal{M}$ ($m = N_{\text{ch}} \times N_{\text{sf}}$ pairs); ToA matrix $P$; Energy matrix $E$; Hypercycle $H$; Gateway demodulation limit $N_{\text{demod}}^{\max} = 8$\\
\textbf{Output:} Feasible schedule $\mathcal{S} = \{(J_i, \alpha(i), s_i, f_i)\}$, Dropped packet set $\mathcal{J}_{\text{drop}}$
\begin{enumerate}\setlength{\itemsep}{1pt}\setlength{\parskip}{0pt}
    \item $\mathcal{J}_{\text{cand}} \leftarrow \mathcal{J}$, $\mathcal{J}_{\text{drop}} \leftarrow \emptyset$.
    \item Clamp deadlines: $D_i \leftarrow \min(D_i, H)$ for all $J_i \in \mathcal{J}_{\text{cand}}$.
    \item \textbf{while} $\mathcal{J}_{\text{cand}} \neq \emptyset$ \textbf{do}
    \item \quad \textbf{// Phase 1: Unrelated Machine Makespan Minimization}
    \item \quad Execute binary search over $w \in [w_{\min}, w_{\max}]$ using $\text{Unsplittable-Truemper}$:
    \item \quad\quad Find assignment $\alpha: \mathcal{J}_{\text{cand}} \to \mathcal{M}$ and critical threshold $w^*$.
    \item \quad\quad Compute machine makespan $M_{\text{span}} \leftarrow \max_{j \in \mathcal{M}} \sum_{i: \alpha(i)=j} p_{ij}$.
    \item \quad \textbf{if} $M_{\text{span}} > H$ \textbf{then}
    \item \quad\quad Identify candidate packet with lowest priority: $J^* \leftarrow \arg\min_{J_i \in \mathcal{J}_{\text{cand}}} (\Pi_i, -D_i, i)$.
    \item \quad\quad $\mathcal{J}_{\text{drop}} \leftarrow \mathcal{J}_{\text{drop}} \cup \{J^*\}$, $\mathcal{J}_{\text{cand}} \leftarrow \mathcal{J}_{\text{cand}} \setminus \{J^*\}$.
    \item \quad\quad \textbf{continue} \Comment{Restart Phase 1 with reduced candidate set}
    \item \quad \textbf{end if}
    \item \quad \textbf{// Phase 1.5: Energy-Optimal Capacity-Constrained Downshifting}
    \item \quad $\alpha \leftarrow \text{Energy-Optimal-Downshift}(\mathcal{J}_{\text{cand}}, \mathcal{M}, P, E, \alpha, w^*, H)$.
    \item \quad \textbf{// Phase 2: Chronological EDF Placement with Segment Tree Demodulation Tracking}
    \item \quad \textbf{while} $\mathcal{J}_{\text{cand}} \neq \emptyset$ \textbf{do}
    \item \quad\quad Initialize lazy Segment Tree over $[0, H]$ and channel cursors $\text{cursor}_j \leftarrow 0$.
    \item \quad\quad Simulate chronological placement across all channels using min-heap priority queue.
    \item \quad\quad Identify deadline violators $\mathcal{V}_{\text{dead}}$ and demodulation peak violators $\mathcal{V}_{\text{demod}}$.
    \item \quad\quad \textbf{if} $\mathcal{V}_{\text{dead}} = \emptyset$ \textbf{and} $\mathcal{V}_{\text{demod}} = \emptyset$ \textbf{then}
    \item \quad\quad\quad Construct schedule $\mathcal{S} \leftarrow \{(J_i, \alpha(i), s_i, f_i) \mid J_i \in \mathcal{J}_{\text{cand}}\}$.
    \item \quad\quad\quad \textbf{return} ($\mathcal{S}, \mathcal{J}_{\text{drop}}$)
    \item \quad\quad \textbf{end if}
    \item \quad\quad $\mathcal{V}_{\text{offending}} \leftarrow \mathcal{V}_{\text{dead}}$ \textbf{if} $\mathcal{V}_{\text{dead}} \neq \emptyset$ \textbf{else} $\mathcal{V}_{\text{demod}}$.
    \item \quad\quad Drop lowest priority offender: $J^* \leftarrow \arg\min_{J_i \in \mathcal{V}_{\text{offending}}} (\Pi_i, -D_i, i)$.
    \item \quad\quad $\mathcal{J}_{\text{drop}} \leftarrow \mathcal{J}_{\text{drop}} \cup \{J^*\}$, $\mathcal{J}_{\text{cand}} \leftarrow \mathcal{J}_{\text{cand}} \setminus \{J^*\}$.
    \item \quad\quad Prune $J^*$ from $\alpha$ and retry Phase 2 directly without re-running Phase 1.
    \item \quad \textbf{end while}
    \item \textbf{end while}
    \item \textbf{return} ($\emptyset, \mathcal{J}$) \Comment{Infeasible instance}
\end{enumerate}
\end{algorithm}

\subsection{Hierarchical Cyclic Executive Architecture and Inter-Cycle Carryover Buffer}
\label{subsec:cyclic_executive}

In real-world LoRaWAN deployments, traffic is inherently periodic and runs indefinitely across repeating major hypercycles $H = \text{lcm}(T_k)$. However, scheduling over a monolithic hypercycle horizon can incur latency accumulation for high-frequency streams and lacks adaptive recovery when transient RF bursts occur. To solve this, our framework implements a hierarchical \textbf{Cyclic Executive Architecture} supporting both uniform (homogeneous) and variable-length (heterogeneous) minor frames:

\begin{enumerate}
    \item \textbf{Macro and Micro Time Horizon Partitioning:} The major hypercycle $H$ is partitioned into $K$ minor operational dispatch frames of duration $\tau_k$ ($k \in \{0, 1, \dots, K-1\}$):
    \begin{itemize}
        \item \textbf{Homogeneous (Uniform) Frames:} All minor cycles share an identical duration $\tau_k = \tau$, where $H = K \cdot \tau$ and $H \pmod \tau = 0$.
        \item \textbf{Heterogeneous (Variable-Length) Frames:} Minor cycle durations may differ ($\tau_0, \tau_1, \dots, \tau_{K-1}$ with $\tau_k > 0$) to align with multi-rate physical sensor activation epochs, subject to the partition invariant:
        \begin{equation}
        \sum_{k=0}^{K-1} \tau_k = H
        \end{equation}
    \end{itemize}
    In both cases, frame boundaries are defined by $t_0 = 0$ and $t_{k+1} = t_k + \tau_k$, establishing mutually disjoint execution windows $[t_k, t_{k+1})$ for all $k \in \{0, \dots, K-1\}$.
    
    \item \textbf{Inter-Cycle Carryover Buffer with Slack Viability Filtering:} Under periodic LoRaWAN traffic, deadlines are governed by the stream period: $D_i = r_i + T_i$. When a sensor stream has period $T_i > \tau_k$ (e.g., $T_i \in \{20, 50, 100\}\text{ s}$ with $\tau_k \in [10, 25]\text{ s}$), a packet dropped in minor frame $k$ due to transient demodulator or channel contention does not immediately constitute a deadline violation.
    
    A dropped packet $J_i$ is eligible for deferral into minor frame $k+1$ if and only if its earliest possible finish time in frame $k+1$ is strictly within its hard deadline:
    \begin{equation}
    t_{k+1} + \min_{j \in \mathcal{M}} p_{ij} \le D_i \quad \text{and} \quad h_i + 1 \le h_{\max} \quad \text{and} \quad t_{k+1} < H
    \end{equation}
    where $t_{k+1}$ is the start time of frame $k+1$, and $h_i$ denotes the number of carryover hops traversed by $J_i$. If this condition is violated, the packet cannot meet its deadline on any machine and is discarded as an unrecoverable hard deadline miss.
    
    \item \textbf{Starvation Prevention via Priority Escalation:} To guarantee timely service and bound scheduling jitter, carried packets receive dynamic priority boosting in subsequent cycles:
    \begin{equation}
    \Pi'_i = \Pi_i + h_i \cdot \Delta_{\Pi}
    \end{equation}
    ensuring that deferred packets take precedence over freshly arriving low-priority traffic.
\end{enumerate}

\begin{theorem}[Preservation of the 2-Approximation Guarantee and RF Invariants under Homogeneous and Heterogeneous Cyclic Scheduling]
\label{thm:cyclic_guarantee}
For any minor frame $k \in \{0, \dots, K-1\}$ of duration $\tau_k$ with candidate packet set $\mathcal{J}^{(k)}$ (comprising new arrivals released in $[t_k, t_{k+1})$ and viable carried-over packets), the cyclic executive scheduler guarantees:
\begin{enumerate}
    \item \textbf{Per-Frame 2-Approximation:} The maximum machine load satisfies:
    \begin{equation}
    C_{\max}^{(k)} \le 2 w_k^* \le 2 \cdot \text{OPT}_k \le \tau_k
    \end{equation}
    where $\text{OPT}_k$ is the optimum makespan for $\mathcal{J}^{(k)}$ on machines $\mathcal{M}$ within frame horizon $\tau_k$.
    \item \textbf{Global RF Invariant Preservation:} Because minor frame timelines are mutually disjoint ($[t_k, t_{k+1}) \cap [t_l, t_{l+1}) = \emptyset$ for $k \ne l$), per-channel non-preemption and gateway demodulation occupancy:
    \begin{equation}
    \Omega(t) \le N_{\text{demod}}^{\max} = 8, \quad \forall t \in [0, H]
    \end{equation}
    hold continuously and unconditionally across the entire global hypercycle.
\end{enumerate}
\end{theorem}

\begin{proof}
In each minor frame $k$, candidate packets $\mathcal{J}^{(k)}$ are assigned via Algorithm~\ref{alg:lorawan_scheduler} over the local horizon $\tau_k$. By Theorem~\ref{thm:overall_2approx}, Phase 1 produces an assignment with load bounded by $2 w_k^* \le 2 \cdot \text{OPT}_k$. By Theorem~\ref{thm:energy_preserves_bound}, Phase 1.5 strictly enforces that no machine load exceeds $2 w_k^*$. Because all transmissions assigned to frame $k$ terminate before $t_{k+1}$, no RF transmission overlaps or demodulation spillover can occur across frame boundaries. Thus, global physical constraints and the 2-approximation makespan bound are preserved across all frames, whether durations $\tau_k$ are uniform or variable.
\end{proof}

\section{A Priori Feasibility and Schedulability Analysis via Upper Bound Duality}
\label{sec:schedulability_certificate}

In real-world mission-critical IoT systems, network operators require formal verification that real-time packet deadlines will be satisfied before network deployment. While empirical simulations provide valuable experimental insights, they cannot guarantee correctness across all operational conditions. To address this, we establish an \textbf{analytical, simulation-free feasibility theorem} by leveraging upper-bound duality between our combinatorial engine and the classical Linear Programming relaxation of Lenstra, Shmoys, and Tardos (LST 1990)~\cite{lenstra1990approximation}.

\subsection{The Lenstra--Shmoys--Tardos (1990) LP Relaxation and Rounding Bound}
Consider the parametric continuous linear program $\text{LP}(T)$ for a guessed makespan $T$ and fleet radio energy budget $\mathcal{E}_{\text{budget}}$:
\begin{align}
\sum_{j \in \mathcal{M}} x_{ij} &= 1 && \forall i \in \mathcal{J} \label{eq:lp_assign} \\
\sum_{i \in \mathcal{J}} x_{ij} \cdot p_{ij} &\le T && \forall j \in \mathcal{M} \label{eq:lp_load} \\
x_{ij} &= 0 && \forall (i, j) \text{ with } p_{ij} > T \text{ or } p_{ij} > D_i \label{eq:lp_prune} \\
\sum_{i \in \mathcal{J}} \sum_{j \in \mathcal{M}} x_{ij} \cdot E_{ij} &\le \mathcal{E}_{\text{budget}} \label{eq:lp_energy} \\
x_{ij} &\ge 0 && \forall i \in \mathcal{J}, j \in \mathcal{M}
\end{align}

\begin{theorem}[LST Bipartite Matching Rounding Bound~\cite{lenstra1990approximation,shmoys1993approximation}]
\label{thm:lst_rounding}
If $\text{LP}(T)$ is feasible, any basic feasible solution (extreme point) $x^*$ can be rounded in polynomial time via maximum bipartite matching on the fractional graph $G_F = (\mathcal{J}_F \cup \mathcal{M}_F, E_F)$ into an integral assignment $\hat{x} \in \{0, 1\}^{n \times m}$ such that:
\begin{equation}
\sum_{i \in \mathcal{J}} \hat{x}_{ij} p_{ij} \le T + \max_{i: x^*_{ij} > 0} p_{ij} \le 2T, \quad \forall j \in \mathcal{M}
\end{equation}
and $\sum_{i,j} \hat{x}_{ij} E_{ij} \le \mathcal{E}_{\text{budget}}$.
\end{theorem}

\begin{proof}
The constraint matrix of $\text{LP}(T)$ possesses $n + m + 1$ functional constraints. Consequently, any basic feasible solution has at most $n + m + 1$ non-zero variables ($x^*_{ij} > 0$). Because each of the $n$ jobs requires $\sum_j x_{ij} = 1$, at most $m + 1$ jobs receive fractional assignments. Standard cycle canceling eliminates cycles in the fractional bipartite graph $G_F$ without increasing machine loads or energy, decomposing $G_F$ into a collection of pseudo-trees. A maximum matching in each tree assigns each fractional job to a distinct machine. Since each machine receives at most one fractional job $i^*$, its load satisfies:
\begin{equation}
\hat{L}_j = \sum_{i: x^*_{ij}=1} p_{ij} + p_{i^* j} \le \sum_{i=1}^n x^*_{ij} p_{ij} + p_{i^* j} \le T + \max_{i: x^*_{ij}>0} p_{ij} \le 2T
\end{equation}
Setting $T = \text{OPT}_{\text{LP}} \le \text{OPT}$ yields an independent integral upper bound $X_{\text{LST}} \le 2 \cdot \text{OPT}$.
\end{proof}

\subsection{Closed-Form Schedulability Certificate Without Simulation}

To enable embedded gateways to verify feasibility in $\mathcal{O}(n \cdot m)$ time without solving an LP, we define closed-form workload parameters:
\begin{itemize}
    \item Minimum ToA per packet: $p_i^{\min} = \min_{j \in \mathcal{M}} p_{ij}$.
    \item Total minimal fleet workload: $W = \sum_{i=1}^n p_i^{\min}$.
    \item Average machine workload: $L_{\text{avg}} = W / m$.
    \item Maximum single-packet transmission time: $p_{\max} = \max_{i=1}^n p_i^{\min}$.
\end{itemize}

\begin{theorem}[A Priori Feasibility Certificate]
\label{thm:apriori_feasibility}
Define the analytical upper bound $X_{\text{tight}}$ as:
\begin{equation}
\label{eq:x_tight}
X_{\text{tight}} = 2 \left( \left\lceil \frac{W}{m} \right\rceil + p_{\max} \right)
\end{equation}
If $X_{\text{tight}} \le H$, then:
\begin{equation}
C_{\max}^{\text{mach}}(\alpha) \le 2 \cdot \text{OPT} \le X_{\text{tight}} \le H
\end{equation}
All packets are mathematically guaranteed to be schedulable within the hypercycle $H$ with zero deadline drops, \textbf{without requiring algorithm simulation or empirical execution}.
\end{theorem}

\begin{proof}
For any feasible schedule, the optimal makespan is lower-bounded by the average load and the longest job: $\text{OPT} \ge \max(\lceil W/m \rceil, p_{\max})$.
Under the worst-case machine dispersion, the binary search upper bound parameter satisfies $w^* \le \lceil W/m \rceil + p_{\max}$.
By Theorem~\ref{thm:overall_2approx}, our combinatorial scheduler produces an assignment $\alpha$ with machine makespan bounded by:
\begin{equation}
C_{\max}^{\text{mach}}(\alpha) \le 2 w^* \le 2 \left( \left\lceil \frac{W}{m} \right\rceil + p_{\max} \right) = X_{\text{tight}}
\end{equation}
If $X_{\text{tight}} \le H$, then by transitivity $C_{\max}^{\text{mach}}(\alpha) \le H$, completing the proof.
\end{proof}

\subsection{Numerical Evaluation on Standard Deployment}
Evaluating Theorem~\ref{thm:apriori_feasibility} on our standard deployment ($n = 234$ packets, $m = 48$ machines, $H = 100{,}000\text{ ms}$, $\tau = 10{,}000\text{ ms}$):
\begin{itemize}
    \item Total minimal fleet workload: $W = 14{,}565\text{ ms}$.
    \item Average workload per machine: $L_{\text{avg}} = 303.44\text{ ms} \implies \lceil W/m \rceil = 304\text{ ms}$.
    \item Maximum single-packet ToA: $p_{\max} = 98\text{ ms}$.
    \item Analytical Lower Bound: $\text{LB} = \max(304, 98) = 304\text{ ms} \le \text{OPT}$.
    \item Proven Upper Bound: $X_{\text{tight}} = 2 \cdot (304 + 98) = \mathbf{804\text{ ms}}$.
\end{itemize}
Because $X_{\text{tight}} = 804\text{ ms} \ll \tau = 10{,}000\text{ ms} \ll H = 100{,}000\text{ ms}$, the workload achieves a \textbf{99.2\% analytical safety margin}. Feasibility is formally proven before scheduling begins.

\section{Time Complexity Analysis of the Complete Framework}
\label{sec:time_complexity}

In this section, we present a comprehensive, rigorous time complexity derivation for the entire LoRaWAN Hypercycle Scheduling framework. We analyze each algorithmic phase individually and derive the overall system asymptotic bound.

\subsection{Phase-by-Phase Computational Complexity}

\subsubsection{Matrix Initialization and Input Generation}
Computing the Time-on-Air matrix $P$ requires evaluating the Semtech SX1276 closed-form airtime equation for each packet $J_i \in \mathcal{J}$ across all $m = N_{\text{ch}} \times N_{\text{sf}}$ communication resources. With $n = |\mathcal{J}|$ packets and $m$ machines, this initial evaluation requires:
\begin{equation}
T_{\text{init}} = \mathcal{O}(n \cdot m)
\end{equation}
For our standard deployment ($n = 234, m = 48$), this step executes in under 2 milliseconds.

\subsubsection{Phase 1: Combinatorial 2-Approximation Makespan Minimization}
Phase 1 searches for the optimal candidate makespan $w$ using binary search, executing Algorithm~\ref{alg:truemper} ($\text{Unsplittable-Truemper}$) at each step:
\begin{enumerate}
    \item \textbf{Binary Search Bounds:} The search interval is bounded by:
    \begin{equation}
    w_{\min} = \max_{i \in \mathcal{J}} \min_{j \in \mathcal{M}} p_{ij}, \quad w_{\max} = 2 \left\lceil \frac{\sum_i \min_j p_{ij}}{m} \right\rceil + 2 \max_{i, j} p_{ij}
    \end{equation}
    The number of binary search iterations is strictly bounded by:
    \begin{equation}
    K_{\text{bs}} = \mathcal{O}\left(\log(w_{\max} - w_{\min})\right)
    \end{equation}
    In typical LoRaWAN deployments, $w_{\max} - w_{\min} \approx 2000\text{ ms}$, yielding $K_{\text{bs}} \le 11$ iterations.

    \item \textbf{Potential Ceiling and Outer Iterations:} The node potentials $\pi(v)$ are bounded above by $\log_b(m)$, where $b = 1 + 1/m$. Because:
    \begin{equation}
    \log_b(m) = \frac{\ln m}{\ln(1 + 1/m)} \le m \ln m
    \end{equation}
    and all edge costs $c_{ij}$ and potential increments are integers ($\ge 1$), the outer `while' loop of Algorithm~\ref{alg:truemper} executes at most:
    \begin{equation}
    K_{\text{outer}} = \mathcal{O}(m \log m) \quad \text{iterations}
    \end{equation}

    \item \textbf{Dijkstra Shortest Path Computation:} In each outer iteration, Dijkstra's algorithm computes shortest paths on the reversed residual graph $G_\alpha(w) = (V, E_\alpha(w))$, where $|V| = n + m$ and $|E_\alpha(w)| \le A + n \le 2A$. Using a binary heap / priority queue implementation:
    \begin{equation}
    T_{\text{dijkstra}} = \mathcal{O}\left(|E| \log |V|\right) = \mathcal{O}\left(A \log(n + m)\right)
    \end{equation}
    (With a Fibonacci heap, this reduces to $\mathcal{O}(A + (n + m) \log(n + m))$.)

    \item \textbf{Inner Loop (Algorithm~\ref{alg:ubf} Unsplittable-Blocking-Flow):}
    By Theorem~\ref{thm:ubf_complexity}, the inner loop performs at most $m$ phases of job pushes. Each phase executes a backward Breadth-First Search (BFS) over the zero-reduced-cost admissible graph $G'_\alpha(w)$ in time proportional to the number of admissible edges, $\mathcal{O}(A)$. Thus, the inner loop terminates in:
    \begin{equation}
    T_{\text{ubf}} = \mathcal{O}(m A)
    \end{equation}
\end{enumerate}

Combining Dijkstra and the inner blocking flow per outer iteration yields $\mathcal{O}(mA + A \log(n+m)) = \mathcal{O}(mA)$. Multiplying by the potential ceiling bound $K_{\text{outer}} = \mathcal{O}(m \log m)$, each call to $\text{Unsplittable-Truemper}$ runs in:
\begin{equation}
T_{\text{truemper}} = \mathcal{O}(m^2 A \log m)
\end{equation}
Incorporating binary search, the total time complexity of Phase 1 is:
\begin{equation}
\label{eq:phase1_complexity}
T_{\text{Phase 1}} = \mathcal{O}\Big(m^2 A \log m \log(w_{\max} - w_{\min})\Big)
\end{equation}
Since $A \le n \cdot m$, for a fixed gateway hardware architecture ($m = 48$ machines), Phase 1 is \textbf{strictly linear in the number of packets $n$}:
\begin{equation}
T_{\text{Phase 1}} = \mathcal{O}\Big(n \cdot m^3 \log m \log(w_{\max} - w_{\min})\Big) \equiv \mathcal{O}(n)
\end{equation}

\subsubsection{Phase 2: Chronological EDF Placement with Segment Tree Tracking}
Phase 2 verifies per-channel non-preemption and enforces global gateway demodulation limits ($N_{\text{demod}}^{\max} \le 8$):
\begin{enumerate}
    \item \textbf{Earliest Deadline First Sorting:} Sorting the assigned packets $\mathcal{J}_j$ on each machine $M_j$ takes:
    \begin{equation}
    T_{\text{sort}} = \sum_{j=1}^m \mathcal{O}(|\mathcal{J}_j| \log |\mathcal{J}_j|) \le \mathcal{O}(n \log n)
    \end{equation}
    \item \textbf{Multi-Channel Progression via Min-Heap:} A min-heap maintains the next earliest ready packet across all $m$ channel queues. Popping and inserting packets requires:
    \begin{equation}
    T_{\text{heap}} = \mathcal{O}(n \log m)
    \end{equation}
    \item \textbf{Demodulation Occupancy via Lazy Segment Tree:} A segment tree over the discretized hypercycle duration $[0, H]$ has depth $\lceil \log_2 H \rceil$. For each scheduled packet:
    \begin{itemize}
        \item Checking the peak concurrent demodulations $\max_{t \in [s_i, f_i-1]} \Omega(t)$ requires one range-maximum query: $\mathcal{O}(\log H)$.
        \item Updating occupancy $\Omega(t) \leftarrow \Omega(t) + 1$ across $[s_i, f_i-1]$ requires one lazy range-addition: $\mathcal{O}(\log H)$.
    \end{itemize}
    Across all $n$ candidate packets, segment tree operations require $\mathcal{O}(n \log H)$.
\end{enumerate}

Thus, a single complete simulation pass of Phase 2 has time complexity:
\begin{equation}
\label{eq:phase2_complexity}
T_{\text{Phase 2}}^{\text{pass}} = \mathcal{O}\Big(n \log n + n \log m + n \log H\Big) = \mathcal{O}\Big(n (\log n + \log m + \log H)\Big)
\end{equation}
With $H = 100{,}000\text{ ms}$, $\log_2 H \approx 17$, making this step practically instantaneous.

\subsubsection{Phase 3: Priority-Conscious Drop Resolution Loop}
When packets violate deadlines or the demodulation threshold, offending packets are pruned:
\begin{enumerate}
    \item \textbf{Drop Selection:} Identifying the lowest-priority offender $J^* = \arg\min (\Pi_i, -D_i, i)$ over the set of violators $\mathcal{V}$ takes linear time: $\mathcal{O}(|\mathcal{V}|) \le \mathcal{O}(n)$.
    \item \textbf{Phase 2 Pruning without Phase 1 Re-run:} In Algorithm~\ref{alg:lorawan_scheduler}, whenever an offending packet is dropped due to timeline or demodulation constraints, it is removed directly from the candidate set $\mathcal{J}_{\text{cand}}$ and pruned from the current assignment $\alpha$. Phase 2 is immediately re-executed on the remaining assignment without re-running the expensive Phase 1.
    \item Across at most $k_{\text{drop}} \le n$ dropped packets, Phase 2 re-simulations consume at most:
    \begin{equation}
    T_{\text{drops}} = \mathcal{O}\Big(k_{\text{drop}} \cdot n (\log m + \log H)\Big) \le \mathcal{O}\Big(n^2 (\log m + \log H)\Big)
    \end{equation}
    \item \textbf{Makespan Infeasibility ($M_{\text{span}} > H$):} If the machine load makespan exceeds the hypercycle duration $H$, Phase 1 must be re-run with one packet removed. In real-time LoRaWAN networks, $M_{\text{span}} \ll H$ (e.g., $1{,}302\text{ ms} \ll 100{,}000\text{ ms}$ in our benchmarks), so Phase 1 is executed \textbf{at most once}.
\end{enumerate}

\subsection{Overall System Time Complexity}

\begin{theorem}[Overall Time Complexity of the Framework]
\label{thm:overall_complexity}
The complete LoRaWAN Hypercycle Scheduling framework (Algorithm~\ref{alg:lorawan_scheduler}) operates within the following asymptotic upper bounds:
\begin{enumerate}
    \item \textbf{Worst-Case Upper Bound:} In the pathological worst case where Phase 1 makespan repeatedly exceeds $H$ triggering $n$ full re-executions:
    \begin{equation}
    \mathcal{O}\Big(n \cdot m^2 A \log m \log(w_{\max} - w_{\min}) + n^2 (\log m + \log H)\Big)
    \end{equation}
    \item \textbf{Standard Operational Bound:} In standard real-time networks where $M_{\text{span}} \le H$ (Phase 1 executes once, and all drops are resolved through direct assignment pruning in Phase 2):
    \begin{equation}
    \label{eq:total_practical_complexity}
    \mathbf{T_{\text{total}}} = \mathcal{O}\Big(m^2 A \log m \log(w_{\max} - w_{\min}) + k_{\text{drop}} \cdot n (\log m + \log H)\Big)
    \end{equation}
    \item \textbf{Hardware-Fixed Asymptotic Bound:} Because the number of gateway channels and spreading factors is fixed by RF hardware ($m = 48$), $m$ is a constant. Expressing $A \le n \cdot m$, the framework scales as:
    \begin{equation}
    \mathbf{T_{\text{total}}} = \mathcal{O}\Big(n \log(w_{\max} - w_{\min}) + k_{\text{drop}} \cdot n \log H\Big) \equiv \mathcal{O}\left(n^2 \log H\right)
    \end{equation}
    When no drops occur ($k_{\text{drop}} = 0$), the entire framework runs in strictly quasi-linear time $\mathcal{O}(n (\log n + \log H))$.
\end{enumerate}
\end{theorem}

\begin{table}[!ht]
\centering
\small
\caption{Detailed Computational Complexity Breakdown by Component}
\label{tab:complexity_breakdown}
\begin{tabular}{|l|l|l|c|}
\hline
\textbf{Component / Phase} & \textbf{Primary Data Structure} & \textbf{Dominant Operation} & \textbf{Asymptotic Time Complexity} \\
\hline
\hline
\textbf{Input Generation} & 2D Float/Int Matrix & SX1276 Airtime Formula & $\mathcal{O}(n \cdot m)$ \\
\hline
\textbf{Dijkstra Shortest Path} & Min-Priority Queue (Heap) & Distance Updates & $\mathcal{O}(A \log(n + m))$ \\
\hline
\textbf{Inner Blocking Flow} & Level Graph + BFS Queue & Admissible Augmentations & $\mathcal{O}(m \cdot A)$ \\
\hline
\textbf{Outer Truemper Loop} & Residual Bipartite Graph & Potential Adjustments & $\mathcal{O}(m^2 A \log m)$ \\
\hline
\textbf{Phase 1 Makespan Opt.} & Binary Search Interval & Unsplittable-Truemper & $\mathcal{O}(m^2 A \log m \log(w_{\max} - w_{\min}))$ \\
\hline
\textbf{Per-Channel EDF Sort} & QuickSort / TimSort & Deadline Ordering & $\mathcal{O}(n \log n)$ \\
\hline
\textbf{Demodulation Tracking} & Lazy Segment Tree & Range-Max / Range-Add & $\mathcal{O}(n \log H)$ \\
\hline
\textbf{Channel Heap Engine} & Min-Heap of Channels & Event Cursor Progression & $\mathcal{O}(n \log m)$ \\
\hline
\textbf{Phase 2 Placement Pass} & Segment Tree + Min-Heap & Chronological Simulation & $\mathcal{O}(n (\log n + \log m + \log H))$ \\
\hline
\textbf{Phase 3 Drop Resolution} & Linear Offender Scan & Prune \& Re-simulate & $\mathcal{O}(k_{\text{drop}} \cdot n (\log m + \log H))$ \\
\hline
\hline
\textbf{Complete System Total} & \textbf{Hybrid Pipeline} & \textbf{End-to-End Execution} & $\mathcal{O}\Big(m^2 A \log m \log \Delta_w + k_{\text{drop}} n \log H\Big)$ \\
\hline
\end{tabular}
\end{table}

\section{Experimental Performance and Empirical Validation}
\label{sec:experiments}

\subsection{Benchmark Configuration and Methodology}
The framework was evaluated on an EU868 LoRaWAN network featuring:
\begin{itemize}
    \item $N_{\text{ch}} = 8$ carrier frequencies ($868.1$ to $868.5\text{ MHz}$, plus 5 additional channels).
    \item $N_{\text{sf}} = 6$ quasi-orthogonal spreading factors ($\text{SF}_7$ to $\text{SF}_{12}$).
    \item Total logical machines: $m = 8 \times 6 = 48$.
    \item Hardware demodulation capacity: $N_{\text{demod}}^{\max} = 8$ concurrent paths.
    \item Hypercycle duration: $H = 100\text{ seconds} = 100{,}000\text{ ms}$.
    \item Traffic workload: 50 independent end-devices generating harmonic traffic with periods $T \in \{10, 20, 50, 100\}\text{ s}$ and payloads $L \in \{10, 20, 30, 50\}\text{ bytes}$, producing $n = 234$ total packet instances.
    \item Deadlines configured with slack factor $1.5 \times \text{ToA}_{\text{SF7}}$.
\end{itemize}

\subsection{Approximation Ratio Verification Against Exact Optimum}
To verify the theoretical 2-approximation bound, an exact Branch-and-Bound (B\&B) mixed-integer solver with aggressive symmetry breaking was executed on the scheduled candidate packet set ($N = 101$).

The exact optimal machine makespan is $\text{OPT} = 800\text{ ms}$. The 2-approximation scheduler achieved a machine load makespan of $1{,}302\text{ ms}$, yielding an empirical ratio of:
\begin{equation}
\frac{C_{\max}^{\text{mach}}(\alpha)}{\text{OPT}} = \frac{1302}{800} = \mathbf{1.6275} \le \mathbf{2.0000}
\end{equation}
The empirical ratio strictly satisfies the theoretical 2.0 bound with a comfortable 18.6\% safety margin. The complete execution completed in $0.412\text{ seconds}$ on standard desktop hardware.

\subsection{Demodulation Occupancy and Invariant Verification}
All system invariants were verified through post-scheduling automated assertion checks:
\begin{enumerate}
    \item \textbf{Channel Non-Preemption:} Pairwise temporal overlaps on the same physical channel and SF evaluated to 0.
    \item \textbf{Deadline Compliance:} Every scheduled packet completed strictly before its deadline: $f_i \le D_i$ for all $101$ scheduled transmissions.
    \item \textbf{Demodulation Capacity:} The Segment Tree tracked continuous gateway occupancy $\Omega(t)$. The peak observed concurrent demodulation across the hypercycle was exactly $8/8$, with zero instances exceeding the threshold:
    \begin{equation}
    \max_{t \in [0, H]} \Omega(t) = 8 \le N_{\text{demod}}^{\max}
    \end{equation}
\end{enumerate}

\subsection{Physical Bottleneck Analysis of Packet Drops}
The simulation scheduled 101 packets and dropped 133 packets (56.8\% drop rate). Rigorous analysis reveals that this is caused by physical hardware saturation rather than algorithmic inefficiency:
\begin{enumerate}
    \item \textbf{Synchronous Release Bottleneck at $t=0$:} All 50 streams release packets simultaneously at $t = 0\text{ ms}$ with tight relative deadlines $D_i \approx 63\text{ ms}$ ($1.5 \times \text{ToA}_{\text{SF7}}$). Because the gateway possesses only 8 physical carrier frequencies, at most 8 packets can physically transmit simultaneously in the interval $[0, 63]\text{ ms}$. Consequently, at least 42 packets at $t=0$ are physically impossible to schedule by \emph{any} algorithm.
    \item \textbf{Spreading Factor Deadline Infeasibility:} Packets with $63\text{ ms}$ relative deadlines cannot utilize SF8 ($74\text{ ms}$), SF9 ($135\text{ ms}$), or SF10--12 because their physical airtime alone exceeds the deadline. Thus, all traffic is restricted to the 8 channels of SF7, saturating the physical gateway capacity.
\end{enumerate}

\subsection{Radio Energy Consumption Evaluation}
\label{subsec:energy_evaluation}
We evaluated the fleet-wide radio energy consumption using the Casals et al.~\cite{casals2017modeling} Class A energy model. Table~\ref{tab:energy_comparison} compares the total energy consumed across all $n = 234$ packets under our 2-approximation framework against the four baseline heuristics.

\begin{table}[!ht]
\centering
\small
\caption{Energy Consumption Comparison across All $n=234$ Packets ($m=48$ Machines)}
\label{tab:energy_comparison}
\begin{tabular}{|l|c|c|c|}
\hline
\textbf{Algorithm} & \textbf{Total Radio Energy (mJ)} & \textbf{Avg Energy (mJ/packet)} & \textbf{Energy Ratio vs. Ours} \\
\hline
\hline
\textbf{2-Approx (Phase 1.5 Ours)} & \textbf{6,493.8} & \textbf{27.8} & \textbf{1.00$\times$ (Baseline)} \\
\hline
Min-Min Heuristic & 8,734.4 & 37.3 & 1.35$\times$ (+34.5\%) \\
\hline
Greedy Least-Loaded (List) & 9,261.2 & 39.6 & 1.43$\times$ (+42.6\%) \\
\hline
Max-Min Heuristic & 9,422.7 & 40.3 & 1.45$\times$ (+45.1\%) \\
\hline
Random Assignment (ALOHA) & 28,838.2 & 123.2 & 4.44$\times$ (+344.1\%) \\
\hline
\end{tabular}
\end{table}

This reveals a critical design insight: while standard heuristics (such as Min-Min and Max-Min) achieve marginally lower load makespan on benign instances by greedily offloading packets onto higher spreading factors (SF8--SF10), they suffer from severe energy inflation (up to $+45.1\%$ more energy consumed). Our framework confines traffic to the lowest feasible spreading factor within the strict $2 w^*$ makespan envelope, achieving the lowest energy footprint among all coordinated algorithms while simultaneously offering a provable worst-case makespan guarantee ($C_{\max} \le 2 \cdot \text{OPT}$). For the scheduled subset ($N = 101$ packets), total radio energy is $2,649.6\text{ mJ}$ ($26.2\text{ mJ/packet}$).

\section{Conclusion and Future Work}
\label{sec:conclusion}

In this report, we presented a comprehensive combinatorial approximation framework for real-time packet scheduling in LoRaWAN networks. By formulating $\langle \text{SF}, \text{CH} \rangle$ pairs as unrelated parallel machines, we implemented Gairing et al.'s 2-approximation algorithm based on Truemper's logarithmic transformation and unsplittable blocking flows. We integrated a lazy-propagated Segment Tree to enforce gateway demodulation ceilings ($\le 8$) in $O(\log H)$ time, and provided a complete step-by-step time complexity derivation establishing that the framework operates in quasi-linear time for fixed gateway hardware.

Empirical validation against an exact Branch-and-Bound solver confirmed that the scheduler strictly respects the $2.0\times$ theoretical guarantee across all instance sizes, achieving an empirical approximation ratio of $1.6275\times$ on full-scale synthetic LoRaWAN workloads ($N = 101$).

Future work will explore:
\begin{enumerate}
    \item Dynamic transmission power control (TPC) coupled with spreading factor selection to co-optimize energy consumption and makespan.
    \item Runtime slack reclamation to accommodate sporadic non-periodic alarm messages.
    \item Weakly-hard real-time guarantees ($(m, k)$-firm deadlines) to bound consecutive packet drops in congested environments.
\end{enumerate}

\bibliographystyle{IEEEtran}
\bibliography{refr}

\end{document}
